\documentclass[11pt]{article}
\usepackage[a4paper,margin=1in]{geometry}
\usepackage{amsmath,amssymb,mathtools,bm}
\usepackage{enumitem,booktabs,array}
\usepackage{tcolorbox}
\usepackage{hyperref}
\usepackage[protrusion=true,expansion=false]{microtype}
\hypersetup{colorlinks=true,linkcolor=blue,urlcolor=blue,citecolor=blue}
\setlist[itemize]{topsep=2pt,itemsep=2pt}
\setlist[enumerate]{topsep=2pt,itemsep=2pt}
\newtcolorbox{keybox}{colback=blue!3!white,colframe=blue!40!black,boxrule=0.4pt,arc=1mm}
\newtcolorbox{alertbox}{colback=red!3!white,colframe=red!40!black,boxrule=0.4pt,arc=1mm}
\newtcolorbox{answerbox}{colback=green!3!white,colframe=green!40!black,boxrule=0.4pt,arc=1mm}
\newtcolorbox{drillbox}{colback=yellow!5!white,colframe=orange!60!black,boxrule=0.4pt,arc=1mm}
\newcommand{\EV}{\mathbb{E}}
\newcommand{\Var}{\mathrm{Var}}
\newcommand{\Cov}{\mathrm{Cov}}
\newcommand{\blank}[1]{\underline{\hspace{#1}}}

\title{E08-03: Andrei, Mann \& Moyen (2019, JFE) \\
\large ``Why Did the $q$ Theory of Investment Start Working?'' --- Active Recall Workbook}
\author{Empirical Corporate Finance II --- Exam Prep}
\date{\today}
\begin{document}\maketitle

\begin{keybox}
\textbf{Paper in one sentence.} AMM show that the investment-$q$ relation has grown stronger over time: in the cross-section the investment-$q$ sensitivity has roughly doubled since the 1970s, and the improvement is driven by a reduction in measurement error in $q$ (better disclosure, less noisy prices, more informed capital markets) rather than structural changes in firm technology.

\textbf{Placement.} Complements Peters-Taylor (E08-02) on intangibles by identifying a \emph{time trend} in investment-$q$ fit; together they explain why modern investment-$q$ regressions work better than early ones.
\end{keybox}

\section*{Round 1: Complete Walk-Through}

\subsection*{1.1 Observation}
Early empirical investment-$q$ regressions (Hayashi 1982, Fazzari et al.\ 1988) had low $R^2$ and modest coefficients. Later work (Peters-Taylor 2017, Eberly et al.\ 2008) shows stronger fit. Is this a real change?

\subsection*{1.2 Hypothesis}
Two candidates:
\begin{itemize}
\item \textbf{Structural}: firms have changed (more intangibles, different technology), altering the true $q$-investment mapping.
\item \textbf{Measurement}: market prices have become more informative about fundamentals, so measured $q$ is a better proxy for marginal Tobin's $q$.
\end{itemize}

\subsection*{1.3 Data and design}
\begin{itemize}
\item Compustat firm-year panel 1950--2015.
\item Standard $q$ and investment measures; subsample analyses by decade.
\item Measures of market informativeness: analyst coverage, disclosure quality, institutional ownership, pricing-efficiency metrics (PIN, price synchronicity).
\item Structural decomposition of investment-$q$ estimates using a measurement-error framework.
\end{itemize}

\subsection*{1.4 Headline results}
\begin{itemize}
\item Investment-$q$ coefficient increases from $\sim$0.02 in 1970s to $\sim$0.05 in 2000s+.
\item $R^2$ roughly doubles across the same period.
\item Trend correlates strongly with indicators of market informativeness.
\item Adding controls for structural changes (industry composition, intangibles) leaves most of the trend unexplained.
\end{itemize}

\subsection*{1.5 Mechanism}
\begin{itemize}
\item Better price efficiency $\Rightarrow$ less measurement error in $q$.
\item With less measurement error, the estimated investment-$q$ coefficient attenuates less toward zero.
\item Investment-$q$ ``puzzle'' was partially a measurement-era artifact.
\end{itemize}

\subsection*{1.6 Implications}
\begin{itemize}
\item Cross-period comparisons of investment sensitivity must account for market informativeness.
\item Research on firm investment should use modern $q$ estimates.
\item Complements Peters-Taylor intangibles story: both measurement and model errors were present earlier.
\end{itemize}

\subsection*{1.7 Caveats}
\begin{alertbox}
\begin{itemize}
\item Informativeness metrics are noisy proxies for true pricing efficiency.
\item Trends in institutional ownership may reflect underlying structural changes.
\item Difficult to cleanly separate measurement and structural effects.
\end{itemize}
\end{alertbox}

\section*{Round 2: Level 1 Fill-in}
\begin{enumerate}
\item Investment-$q$ coefficient trend: $0.02\to$\blank{1cm} over decades.
\item Main explanation: \blank{2cm} error reduction in $q$.
\item Proxies for market \blank{2cm}: analyst coverage, PIN, disclosure.
\item Sample: Compustat \blank{1cm}--\blank{1cm}.
\item Complement paper: Peters-\blank{1cm} 2017.
\end{enumerate}

\section*{Round 3: Level 2 Prompts}
\begin{enumerate}
\item Describe the observed trend in investment-$q$ fit.
\item State the measurement-vs-structural hypotheses.
\item Report main results distinguishing the two.
\item Discuss implications for applied investment research.
\item Relate to Peters-Taylor.
\end{enumerate}

\section*{Round 4: Minimal Scaffolding}
From a blank page: (i) puzzle, (ii) hypotheses, (iii) data, (iv) results, (v) implications.

\section*{Round 5: Blank Page}
Essay: why has $q$-theory started working? Measurement vs.\ structure.

\vspace{6pt}
\begin{drillbox}
\textbf{Speed Drill.} (1) Coefficient trend. (2) Two hypotheses. (3) Key proxies. (4) Main conclusion. (5) Complement paper.
\end{drillbox}

\section*{Quick Reference \& Answer Key}
\begin{answerbox}
\begin{itemize}
\item \textbf{Trend.} Investment-$q$ slope $\sim$ doubles.
\item \textbf{Cause.} Reduced $q$ measurement error.
\item \textbf{Proxy.} Market informativeness rises.
\item \textbf{Lesson.} Not structural change, mostly measurement.
\end{itemize}
\end{answerbox}

\begin{keybox}
\textbf{30-second exam pitch.} Andrei, Mann \& Moyen (2019) document that the empirical investment-$q$ relation has strengthened over the past half-century --- coefficient roughly doubled, $R^2$ roughly doubled --- and ask whether this reflects structural changes in firms or improvement in measurement. Using Compustat 1950--2015 and proxies for market price informativeness (analyst coverage, institutional ownership, disclosure quality, pricing-efficiency metrics), they show the trend tracks market-informativeness gains rather than shifts in industry mix or intangibles alone. Under classical measurement-error attenuation, better-measured $q$ produces estimates closer to the true coefficient. The paper reframes the investment-$q$ puzzle as partially an era-specific measurement artifact, and complements Peters-Taylor (2017) on intangibles: both measurement (AMM) and model misspecification (PT) contributed to the original puzzle, and modern data + methods resolve both.
\end{keybox}

\end{document}
